\section{总结与延伸}

\subsection{核心要点回顾}

\begin{itemize}[leftmargin=2em]
    \item 记忆是让 AI Agent 从"工具"升级为"伙伴"的关键能力
    \item Agentic Memory 同时服务于三个目标：连续性、上下文和学习
    \item 四种记忆类型（In-context、External、Episodic、Parametric）各司其职，协同工作
    \item 记忆架构的核心在于\textbf{检索设计}而非存储
    \item 记忆管理（遗忘策略）与记忆存储同等重要
\end{itemize}

\subsection{延伸思考}

\begin{itemize}[leftmargin=2em]
    \item \textbf{MemGPT}：将操作系统的虚拟内存概念引入 LLM Agent，自动在不同层级的记忆之间进行 paging
    \item \textbf{Generative Agents}（Stanford）：在模拟世界中使用 Episodic Memory + Reflection 机制让 Agent 展现出类人的社会行为
    \item \textbf{RAG 与记忆的关系}：Retrieval-Augmented Generation 可以看作 External Memory 的一种应用形态，但完整的 Agentic Memory 还包括写入、遗忘和反思机制
    \item \textbf{Multi-Agent 场景下的共享记忆}：当多个 Agent 协作时，如何设计共享的记忆空间和权限控制是一个开放问题
\end{itemize}

%% --- Fin del contenido --- %%

\end{document}
